% ============================================================
% 第8章补充：经典例题与自学元素
% 插入位置：第8章现有miniquiz之后（\end{miniquiz} 之后）
% ============================================================

\section*{习题精选与详解}
\addcontentsline{toc}{section}{习题精选}

\begin{example}{库仑对数在不同等离子体环境中的计算}{coulomb-log-calc}
分别计算以下三种等离子体的库仑对数 $\ln\Lambda$：
\begin{itemize}[leftmargin=*,itemsep=0pt]
\item 实验室低温等离子体：$n_e = 10^{18}\,\mathrm{m}^{-3}$，$T_e = 4\,\mathrm{eV}$
\item 太阳日冕：$n_e = 10^{12}\,\mathrm{m}^{-3}$，$T_e = 100\,\mathrm{eV}$
\item 托卡马克聚变等离子体：$n_e = 10^{20}\,\mathrm{m}^{-3}$，$T_e = 10\,\mathrm{keV}$
\end{itemize}
并讨论 $\ln\Lambda$ 随等离子体参数的变化规律。
\tcblower
\textbf{解答：}

库仑对数定义为 $\ln\Lambda = \ln(12\pi N_D) = \ln(12\pi n_e \lambda_D^3)$，其中 $N_D = n_e \lambda_D^3$ 是德拜球内粒子数。另一种常用形式是
\[
\Lambda = \frac{\lambda_D}{b_0} = \frac{12\pi\eps_0 k_{\!B}T_e}{e^2 n_e^{1/3}}
\]

对典型等离子体，$\ln\Lambda$ 在 $10$--$30$ 之间。下面用 $N_D = n_e \lambda_D^3$ 计算。

已知 $\lambda_D = \sqrt{\frac{\eps_0 k_{\!B}T_e}{n_e e^2}}$，故
\[
N_D = n_e \lambda_D^3 = n_e \left(\frac{\eps_0 k_{\!B}T_e}{n_e e^2}\right)^{3/2}
= \frac{(\eps_0 k_{\!B}T_e)^{3/2}}{n_e^{1/2}e^3}
\]

\textbf{（1）实验室低温等离子体：}
\[
N_D = 10^{18}\times\left(\sqrt{\frac{8.854\times 10^{-12}\times 6.408\times 10^{-19}}{10^{18}\times 2.566\times 10^{-38}}}\right)^3
= 10^{18}\times (1.49\times 10^{-5})^3
\approx 3.3\times 10^3
\]
\[
\ln\Lambda = \ln(12\pi\times 3.3\times 10^3) = \ln(1.24\times 10^5) \approx 11.7
\]

\textbf{（2）太阳日冕：}
$k_{\!B}T_e = 100\,\mathrm{eV} = 1.602\times 10^{-17}\,\mathrm{J}$
\[
\lambda_D = \sqrt{\frac{8.854\times 10^{-12}\times 1.602\times 10^{-17}}{10^{12}\times 2.566\times 10^{-38}}}
= \sqrt{5.53\times 10^{-3}} \approx 7.43\times 10^{-2}\,\mathrm{m} = 7.43\,\mathrm{cm}
\]
\[
N_D = 10^{12}\times (7.43\times 10^{-2})^3 \approx 10^{12}\times 4.1\times 10^{-4} \approx 4.1\times 10^8
\]
\[
\ln\Lambda = \ln(12\pi\times 4.1\times 10^8) = \ln(1.55\times 10^{10}) \approx 23.5
\]

\textbf{（3）托卡马克聚变等离子体：}
$k_{\!B}T_e = 10^4\,\mathrm{eV} = 1.602\times 10^{-15}\,\mathrm{J}$
\[
\lambda_D = \sqrt{\frac{8.854\times 10^{-12}\times 1.602\times 10^{-15}}{10^{20}\times 2.566\times 10^{-38}}}
\approx 7.43\times 10^{-5}\,\mathrm{m} = 74.3\,\mu\mathrm{m}
\]
\[
N_D = 10^{20}\times (7.43\times 10^{-5})^3 \approx 10^{20}\times 4.1\times 10^{-13} \approx 4.1\times 10^7
\]
\[
\ln\Lambda = \ln(12\pi\times 4.1\times 10^7) = \ln(1.55\times 10^9) \approx 21.2
\]

\textbf{规律总结：}
\begin{itemize}[leftmargin=*,itemsep=0pt]
\item $\ln\Lambda$ 对等离子体参数的变化极为缓慢：跨越十几个数量级的密度和温度，$\ln\Lambda$ 仅从约 $12$ 变到约 $23$。
\item 高温低密度等离子体（日冕）$\ln\Lambda$ 较大，因为德拜球内粒子数更多。
\item 实验室等离子体 $\ln\Lambda$ 较小，但仍 $\gg 1$，保证小角度散射主导、两体碰撞近似有效。
\item 在大多数计算中，$\ln\Lambda$ 可以取为 $15$--$20$ 的常数，误差不大。
\end{itemize}
\end{example}

\begin{example}{Spitzer电阻率的数值计算与对比}{spitzer-calc}
某托卡马克装置中，等离子体参数为：$n_e = 5\times 10^{19}\,\mathrm{m}^{-3}$，$T_e = 2\,\mathrm{keV}$，$Z_{\text{eff}} = 1.5$（有效电荷数），库仑对数 $\ln\Lambda = 18$。请计算：
\begin{enumerate}[label=(\arabic*)]
\item 该等离子体的Spitzer电阻率 $\eta$；
\item 对应的电导率 $\sigma = 1/\eta$；
\item 与铜（$\eta_{\text{Cu}} = 1.7\times 10^{-8}\,\Omega\cdot\mathrm{m}$）和海水（$\eta_{\text{seawater}} \approx 0.2\,\Omega\cdot\mathrm{m}$）的电阻率进行对比。
\end{enumerate}
\tcblower
\textbf{解答：}

Spitzer电阻率的数值公式（平行方向）为：
\[
\eta_{\parallel} \approx 5.2\times 10^{-5}\frac{Z\ln\Lambda}{T_e^{3/2}\,\mathrm{(eV)}}\,\Omega\cdot\mathrm{m}
\]

\textbf{(1) 计算Spitzer电阻率：}

将 $T_e = 2\,\mathrm{keV} = 2000\,\mathrm{eV}$，$Z = 1.5$，$\ln\Lambda = 18$ 代入：
\[
\eta = 5.2\times 10^{-5}\times\frac{1.5\times 18}{(2000)^{3/2}}
\]

计算分母：
\[
(2000)^{3/2} = 2000\times\sqrt{2000} = 2000\times 44.72 \approx 8.94\times 10^4
\]

计算分子：
\[
1.5\times 18 = 27
\]

\[
\eta = 5.2\times 10^{-5}\times\frac{27}{8.94\times 10^4}
= 5.2\times 10^{-5}\times 3.02\times 10^{-4}
\approx 1.57\times 10^{-8}\,\Omega\cdot\mathrm{m}
\]

\textbf{(2) 计算电导率：}
\[
\sigma = \frac{1}{\eta} \approx \frac{1}{1.57\times 10^{-8}} \approx 6.4\times 10^{7}\,\mathrm{S/m}
\]

\textbf{(3) 对比：}
\begin{itemize}[leftmargin=*,itemsep=0pt]
\item 铜：$\eta_{\text{Cu}} = 1.7\times 10^{-8}\,\Omega\cdot\mathrm{m}$
\item 该等离子体：$\eta \approx 1.57\times 10^{-8}\,\Omega\cdot\mathrm{m}$
\item 海水：$\eta_{\text{seawater}} \approx 0.2\,\Omega\cdot\mathrm{m}$
\end{itemize}

\textbf{分析：}
\begin{itemize}[leftmargin=*,itemsep=0pt]
\item 在 $T_e = 2\,\mathrm{keV}$ 时，聚变等离子体的电阻率已与铜相当！这说明Spitzer电阻率随温度下降极快：$\eta \propto T_e^{-3/2}$。
\item 若温度升高到 $10\,\mathrm{keV}$，电阻率将下降为 $(10/2)^{3/2} = 11.2$ 倍，即约 $1.4\times 10^{-9}\,\Omega\cdot\mathrm{m}$，仅为铜的约 $1/12$。
\item 海水的电阻率比等离子体高约 $10^{10}$ 倍，这说明即使是"高温"等离子体，其导电能力也远非日常生活中的盐水可比。
\item 电阻率与密度无关：若将密度加倍，载流子数加倍，但碰撞频率也加倍，两者抵消。
\end{itemize}
\end{example}

\begin{example}{双极扩散系数的计算}{ambipolar-calc}
一弱电离氩等离子体中，电子温度 $T_e = 3\,\mathrm{eV}$，离子温度 $T_i = 0.3\,\mathrm{eV}$（即 $T_e/T_i = 10$），电子密度 $n_e = 10^{18}\,\mathrm{m}^{-3}$。已知电子扩散系数 $D_e = 0.5\,\mathrm{m}^2/\mathrm{s}$，离子扩散系数 $D_i = 2.5\times 10^{-3}\,\mathrm{m}^2/\mathrm{s}$。求：
\begin{enumerate}[label=(\arabic*)]
\item 双极扩散系数 $D_a$；
\item 电子和离子在无耦合情况下各自扩散 $1\,\mathrm{cm}$ 所需的时间；
\item 双极扩散下共同扩散 $1\,\mathrm{cm}$ 所需的时间。
\end{enumerate}
\tcblower
\textbf{解答：}

\textbf{(1) 计算双极扩散系数：}

双极扩散系数公式为：
\[
D_a = \frac{\mu_i D_e + \mu_e D_i}{\mu_i + \mu_e}
\]

在 $T_e \gg T_i$ 时，可近似为：
\[
D_a \approx D_i\left(1 + \frac{T_e}{T_i}\right)
\]

代入 $D_i = 2.5\times 10^{-3}\,\mathrm{m}^2/\mathrm{s}$，$T_e/T_i = 10$：
\[
D_a = 2.5\times 10^{-3}\times (1 + 10) = 2.5\times 10^{-3}\times 11 = 2.75\times 10^{-2}\,\mathrm{m}^2/\mathrm{s}
\]

对比：$D_a \approx 11 D_i \ll D_e$。电子扩散被离子"拖慢"，离子扩散被电子温度"加速"。

\textbf{(2) 无耦合扩散时间：}

扩散的特征时间尺度：$\tau \sim L^2/D$。取 $L = 1\,\mathrm{cm} = 0.01\,\mathrm{m}$：

电子：
\[
\tau_e = \frac{L^2}{D_e} = \frac{(0.01)^2}{0.5} = \frac{10^{-4}}{0.5} = 2\times 10^{-4}\,\mathrm{s} = 0.2\,\mathrm{ms}
\]

离子：
\[
\tau_i = \frac{L^2}{D_i} = \frac{10^{-4}}{2.5\times 10^{-3}} = 0.04\,\mathrm{s} = 40\,\mathrm{ms}
\]

\textbf{(3) 双极扩散时间：}
\[
\tau_a = \frac{L^2}{D_a} = \frac{10^{-4}}{2.75\times 10^{-2}} \approx 3.64\times 10^{-3}\,\mathrm{s} \approx 3.6\,\mathrm{ms}
\]

\textbf{物理分析：}
\begin{itemize}[leftmargin=*,itemsep=0pt]
\item 无耦合时，电子只需 $0.2\,\mathrm{ms}$ 就能扩散 $1\,\mathrm{cm}$，而离子需要 $40\,\mathrm{ms}$。
\item 双极扩散下，电子和离子以相同速率 $D_a$ 扩散，耗时约 $3.6\,\mathrm{ms}$。这个时间介于两者之间，但更接近离子的扩散时间（因为离子是"瓶颈"）。
\item 电子温度越高（$T_e/T_i$ 越大），双极扩散系数越大，等离子体边界损失越快。这是放电等离子体（如霓虹灯、ICP刻蚀机）设计中的关键参数。
\end{itemize}
\end{example}

\begin{example}{电子--离子碰撞频率随温度的变化}{collision-freq-temp}
某氢等离子体中，电子密度 $n_e = 10^{20}\,\mathrm{m}^{-3}$，库仑对数 $\ln\Lambda = 17$。计算电子--离子碰撞频率 $\nu_{ei}$ 在 $T_e = 100\,\mathrm{eV}$、$1\,\mathrm{keV}$ 和 $10\,\mathrm{keV}$ 时的值，并分析其温度依赖关系。
\tcblower
\textbf{解答：}

电子--离子碰撞频率公式为：
\[
\nu_{ei} = \frac{n_i Z^2 e^4 \ln\Lambda}{12\pi^{3/2}\eps_0^2 m_e^{1/2}(k_{\!B}T_e)^{3/2}}
\]

对氢等离子体 $Z=1$，$n_i = n_e = 10^{20}\,\mathrm{m}^{-3}$。整理为数值形式：
\[
\nu_{ei} \approx 2.9\times 10^{-12}\,\frac{n_e\ln\Lambda}{T_e^{3/2}\,\mathrm{(eV)}}\,\mathrm{s}^{-1}
\]

\textbf{(1) $T_e = 100\,\mathrm{eV}$：}
\[
\begin{aligned}
\nu_{ei} &= 2.9\times 10^{-12}\times\frac{10^{20}\times 17}{(100)^{3/2}}
= 2.9\times 10^{-12}\times\frac{1.7\times 10^{21}}{10^3}\\
&= 2.9\times 10^{-12}\times 1.7\times 10^{18}
\approx 4.9\times 10^{6}\,\mathrm{s}^{-1}
\end{aligned}
\]

对应碰撞时间：$\tau_{ei} = 1/\nu_{ei} \approx 2.0\times 10^{-7}\,\mathrm{s} = 0.2\,\mu\mathrm{s}$。

\textbf{(2) $T_e = 1\,\mathrm{keV} = 1000\,\mathrm{eV}$：}
\[
\nu_{ei} = 2.9\times 10^{-12}\times\frac{1.7\times 10^{21}}{(1000)^{3/2}}
= 2.9\times 10^{-12}\times\frac{1.7\times 10^{21}}{3.16\times 10^4}
\approx 1.56\times 10^{5}\,\mathrm{s}^{-1}
\]

对应 $\tau_{ei} \approx 6.4\,\mu\mathrm{s}$。

\textbf{(3) $T_e = 10\,\mathrm{keV} = 10^4\,\mathrm{eV}$：}
\[
\nu_{ei} = 2.9\times 10^{-12}\times\frac{1.7\times 10^{21}}{(10^4)^{3/2}}
= 2.9\times 10^{-12}\times\frac{1.7\times 10^{21}}{10^6}
\approx 4.9\times 10^{3}\,\mathrm{s}^{-1}
\]

对应 $\tau_{ei} \approx 2.0\times 10^{-4}\,\mathrm{s} = 0.2\,\mathrm{ms}$。

\textbf{总结：}
\[
\begin{array}{c|c|c}
T_e & \nu_{ei} & \tau_{ei} \\ \hline
100\,\mathrm{eV} & 4.9\times 10^{6}\,\mathrm{s}^{-1} & 0.2\,\mu\mathrm{s} \\
1\,\mathrm{keV} & 1.6\times 10^{5}\,\mathrm{s}^{-1} & 6.4\,\mu\mathrm{s} \\
10\,\mathrm{keV} & 4.9\times 10^{3}\,\mathrm{s}^{-1} & 0.2\,\mathrm{ms}
\end{array}
\]

\textbf{物理分析：} $\nu_{ei} \propto T_e^{-3/2}$，温度升高一个数量级，碰撞频率下降约 $30$ 倍。这是因为高温电子速度更快，在靶离子附近停留时间更短，偏转更小。在聚变等离子体（$T_e\sim 10\,\mathrm{keV}$）中，经典碰撞极为稀少，碰撞输运很弱——这解释了为什么等离子体是近乎理想的导体，也说明了为什么反常输运（由不稳定性导致）成为主要损失机制。
\end{example}

\begin{figure}[htbp]
\centering
\begin{tikzpicture}[scale=1.1,font=\small]
% 入射粒子轨迹
\draw[->,thick,blue] (-4,2) -- (-1.5,0.5);
\draw[->,thick,blue] (-1.5,0.5) -- (3,-1);
\draw[->,thick,blue] (-4,-2) -- (-1.5,-0.5);
\draw[->,thick,blue] (-1.5,-0.5) -- (3,1);

% 靶粒子
\fill[red] (0,0) circle (4pt);
\node[red,right] at (0.3,0) {$q_2$};

% 入射粒子标注
\node[blue,left] at (-4,2.2) {$q_1$, $v$};
\node[blue,left] at (-4,-2.2) {$q_1$, $v$};

% 碰撞参数
\draw[dashed,gray] (-1.5,0.5) -- (-1.5,-0.5);
\node[gray,right] at (-1.3,0) {$b$};

% 偏转角标注
\draw[dashed] (0,0) -- (-1,0.67);
\draw[dashed] (0,0) -- (2,-0.67);
\draw[->,thick,red] (0.8,0) arc (0:-18:0.8);
\node[red] at (1.2,-0.25) {$\theta$};

% 另一偏转角
\draw[dashed] (0,0) -- (-1,-0.67);
\draw[dashed] (0,0) -- (2,0.67);
\draw[->,thick,red] (0.8,0) arc (0:18:0.8);
\node[red] at (1.2,0.25) {$\theta$};

% 小角度示意
\node[align=center] at (0,-2.5) {\footnotesize 单次散射偏转角很小：$\theta \ll 1$};
\node[align=center] at (0,-3) {\footnotesize 多次小角度散射累积等效于一次大角度偏转};
\end{tikzpicture}
\caption{库仑碰撞小角度散射示意图。带电粒子 $q_1$ 以速度 $v$ 经过靶粒子 $q_2$ 附近，由于长程库仑力，轨迹发生轻微偏折。偏转角 $\theta$ 很小，但远距离碰撞频繁发生，累积效应主导了动量输运。}
\end{figure}

% ch8 新例题：Spitzer 电阻率随温度的标度与等离子体加热
\begin{example}{欧姆加热与温度的稳态}{ohmic-heating}
实验室氩等离子体由欧姆加热维持：电流密度 $j = 10^4\,\mathrm{A/m^2}$，密度 $n_e = 10^{19}\,\mathrm{m^{-3}}$，库仑对数 $\ln\Lambda = 10$。
（1）写出欧姆加热功率密度 $P = \eta j^2$ 关于 $T_e$ 的表达式，说明其温度标度；
（2）计算 $T_e = 5\,\mathrm{eV}$ 时的欧姆加热功率密度；
（3）解释为什么纯欧姆加热难以把等离子体加热到高温（远高于 $\sim100\,\mathrm{eV}$）。
\tcblower
\textbf{解答：}

\textbf{(1) 温度标度：} Spitzer 电阻率 $\eta \approx 5\times10^{-5}\,Z\ln\Lambda/T_e^{3/2}\,\Omega\mathrm{m}$（$T_e$ 以 eV 计），故
\[
P = \eta j^2 \approx \frac{5\times10^{-5}\,Z\ln\Lambda\,j^2}{T_e^{3/2}} \propto T_e^{-3/2}
\]
欧姆加热功率随温度升高而\textbf{下降}（电阻率随温度下降）。

\textbf{(2) $T_e = 5\,\mathrm{eV}$ 时：}
\[
\eta \approx \frac{5\times10^{-5}\times1\times10}{5^{3/2}} = \frac{5\times10^{-4}}{11.2} \approx 4.5\times10^{-5}\,\Omega\mathrm{m}
\]
\[
P = \eta j^2 = 4.5\times10^{-5}\times(10^4)^2 = 4.5\times10^3\,\mathrm{W/m^3}
\]

\textbf{(3) 高温极限：} 欧姆加热 $\propto T_e^{-3/2}$ 越来越弱，而辐射损失（韧致辐射 $\propto T_e^{1/2}n^2$）和输运损失随温度上升。两者在 $T_e\sim100\,\mathrm{eV}$ 量级达到平衡后，欧姆加热无法进一步提高温度——这就是托卡马克必须采用中性束注入、ICRF/ECD 射频加热等辅助加热手段的根本原因（欧姆加热只能把等离子体带入"温区"，点火必须靠辅助加热）。
\end{example}

\begin{formulabox}{第8章核心公式速查}{ch8-formulas}
\begin{itemize}[leftmargin=*,itemsep=0pt]
\item Rutherford散射角：$\displaystyle\tan\frac{\theta}{2} = \frac{b_0}{2b}$，其中 $b_0 = \frac{|q_1 q_2|}{4\pi\eps_0 m_r v^2}$
\item 电子--离子碰撞频率：$\displaystyle\nu_{ei} = \frac{n_i Z^2 e^4 \ln\Lambda}{12\pi^{3/2}\eps_0^2 m_e^{1/2}(k_{\!B}T_e)^{3/2}}$
\item 库仑对数：$\displaystyle\Lambda = 12\pi n_e\lambda_D^3 = \frac{12\pi\eps_0 k_{\!B}T_e}{e^2 n_e^{1/3}}$
\item Spitzer电阻率：$\displaystyle\eta_{\parallel} \approx 5.2\times 10^{-5}\frac{Z\ln\Lambda}{T_e^{3/2}\,\mathrm{(eV)}}\,\Omega\cdot\mathrm{m}$
\item 双极扩散系数：$\displaystyle D_a = \frac{\mu_i D_e + \mu_e D_i}{\mu_i + \mu_e} \approx D_i\left(1 + \frac{T_e}{T_i}\right)$
\item 扩散系数（Fick定律）：$\displaystyle\Gamma = -D\grad n$
\item 迁移率：$\displaystyle\mu = \frac{q}{m\nu}$
\item 电阻率与温度关系：$\displaystyle\eta \propto T_e^{-3/2}$（与密度无关）
\end{itemize}
\end{formulabox}

\begin{checklistbox}{第8章自学检查清单}{ch8-checklist}
\begin{itemize}[leftmargin=*,label=$\square$]
\item 我能理解库仑碰撞与硬球碰撞的区别：长程力导致小角度散射主导
\item 我能计算给定等离子体参数下的库仑对数 $\ln\Lambda$
\item 我知道为什么 $\ln\Lambda$ 在很宽的参数范围内变化缓慢（$10$--$30$）
\item 我能从Spitzer公式估算等离子体的电阻率，并理解它与密度无关的原因
\item 我能解释双极扩散的物理图像：电子扩散快 $\to$ 电荷分离 $\to$ 电场拉回 $\to$ 共同扩散
\item 我知道碰撞频率 $\nu_{ei} \propto T_e^{-3/2}$，并理解高温等离子体为何近乎理想导体
\item 我能区分经典输运（碰撞主导）和反常输运（不稳定性/湍流主导）
\end{itemize}
\end{checklistbox}

\endinput
